% ============================================================
%  Praeambel + farbige Boxen + C++-Listings
%  fuer den C++-Lernzettel (Programmieren II)
% ============================================================
\usepackage[utf8]{inputenc}
\usepackage[T1]{fontenc}
\usepackage{lmodern}
\usepackage[ngerman]{babel}
\usepackage{microtype}
\usepackage[a4paper,margin=2.2cm]{geometry}

\usepackage{graphicx}
\usepackage{booktabs}
\usepackage{array}
\usepackage{multirow}
\usepackage{tabularx}
\usepackage{amsmath}
\usepackage{amssymb}
\usepackage{enumitem}
\usepackage{xcolor}
\usepackage{listings}
\usepackage{multicol}

% Bildersuchpfad
\graphicspath{{bilder/}}

% Listen etwas kompakter
\setlist{itemsep=2pt,topsep=3pt,parsep=0pt}

% ---- Farben -------------------------------------------------
\definecolor{defblue}{HTML}{1F6FB2}
\definecolor{defbluebg}{HTML}{EAF2FB}
\definecolor{merkered}{HTML}{C0392B}
\definecolor{merkeredbg}{HTML}{FBEAE8}
\definecolor{syntaxgreen}{HTML}{1E8449}
\definecolor{syntaxgreenbg}{HTML}{E8F6EE}
\definecolor{beispielorange}{HTML}{B9770E}
\definecolor{beispielorangebg}{HTML}{FCF3E2}
\definecolor{tippteal}{HTML}{0E7C7B}
\definecolor{ticktealbg}{HTML}{E4F4F4}
\definecolor{rulegray}{HTML}{888888}

% ---- Code-Farben (Syntax-Highlighting) ----------------------
\definecolor{codebg}{HTML}{F7F8FA}
\definecolor{codekw}{HTML}{0033B3}   % keywords
\definecolor{codekw2}{HTML}{7A3E9D}  % typen / sekundaere keywords
\definecolor{codecomment}{HTML}{4C8A4C}
\definecolor{codestring}{HTML}{A31515}
\definecolor{codenum}{HTML}{098658}
\definecolor{codeline}{HTML}{B0B0B0}

% ---- hyperref (vor tcolorbox) -------------------------------
\usepackage[colorlinks=true,linkcolor=defblue,urlcolor=defblue,
            citecolor=defblue,linktoc=all,
            pdftitle={C++ -- Programmieren II -- Lernzettel},
            pdfauthor={Programmieren II (TSA/TSL 25)}]{hyperref}

% ---- tcolorbox + Box-Umgebungen -----------------------------
\usepackage[most]{tcolorbox}
\tcbuselibrary{skins,breakable}

\newtcolorbox{definition}[1][]{
  enhanced, breakable, colback=defbluebg, colframe=defblue,
  coltitle=white, fonttitle=\bfseries, left=2mm, right=2mm,
  title={Definition#1}, sharp corners=downhill, boxrule=0.6pt,
  attach boxed title to top left={xshift=4mm,yshift=-2.5mm},
  boxed title style={colback=defblue,sharp corners}}

\newtcolorbox{merke}[1][]{
  enhanced, breakable, colback=merkeredbg, colframe=merkered,
  coltitle=white, fonttitle=\bfseries, left=2mm, right=2mm,
  title={Merke \textperiodcentered{} Pr\"ufungstipp#1}, sharp corners=downhill, boxrule=0.6pt,
  attach boxed title to top left={xshift=4mm,yshift=-2.5mm},
  boxed title style={colback=merkered,sharp corners}}

\newtcolorbox{syntaxbox}[1][]{
  enhanced, breakable, colback=syntaxgreenbg, colframe=syntaxgreen,
  coltitle=white, fonttitle=\bfseries, left=2mm, right=2mm,
  title={Syntax#1}, sharp corners=downhill, boxrule=0.6pt,
  attach boxed title to top left={xshift=4mm,yshift=-2.5mm},
  boxed title style={colback=syntaxgreen,sharp corners}}

\newtcolorbox{beispiel}[1][]{
  enhanced, breakable, colback=beispielorangebg, colframe=beispielorange,
  coltitle=white, fonttitle=\bfseries, left=2mm, right=2mm,
  title={Beispiel#1}, sharp corners=downhill, boxrule=0.6pt,
  attach boxed title to top left={xshift=4mm,yshift=-2.5mm},
  boxed title style={colback=beispielorange,sharp corners}}

\newtcolorbox{tipp}[1][]{
  enhanced, breakable, colback=ticktealbg, colframe=tippteal,
  coltitle=white, fonttitle=\bfseries, left=2mm, right=2mm,
  title={Best Practice \textperiodcentered{} Konvention#1}, sharp corners=downhill, boxrule=0.6pt,
  attach boxed title to top left={xshift=4mm,yshift=-2.5mm},
  boxed title style={colback=tippteal,sharp corners}}

% ---- Code-Listings (C++) ------------------------------------
\lstdefinestyle{cpp}{
  language=C++,
  basicstyle=\ttfamily\footnotesize,
  keywordstyle=\color{codekw}\bfseries,
  keywordstyle=[2]\color{codekw2},
  commentstyle=\color{codecomment}\itshape,
  stringstyle=\color{codestring},
  numberstyle=\tiny\color{codeline},
  backgroundcolor=\color{codebg},
  frame=single, rulecolor=\color{rulegray},
  breaklines=true, columns=fullflexible,
  keepspaces=true, showstringspaces=false,
  numbers=none, xleftmargin=3mm, xrightmargin=2mm,
  aboveskip=5pt, belowskip=5pt, tabsize=2,
  morekeywords={override,final,nullptr,constexpr,noexcept,explicit,
                using,namespace,template,typename,virtual,public,private,
                protected,friend,this,new,delete,throw,try,catch,auto,
                static_cast,dynamic_cast,const_cast,reinterpret_cast},
  keywords=[2]{std,string,vector,map,unique_ptr,shared_ptr,make_unique,
               make_shared,size_t,uint8_t,cout,cin,endl,ostream,istream,
               runtime_error,out_of_range,invalid_argument,move,
               initializer_list,pair},
  literate={ä}{{\"a}}1 {ö}{{\"o}}1 {ü}{{\"u}}1 {Ä}{{\"A}}1
           {Ö}{{\"O}}1 {Ü}{{\"U}}1 {ß}{{\ss}}1 {°}{{$^\circ$}}1
           {²}{{$^2$}}1}
\lstset{style=cpp}

% Inline-Code als kleines Makro
\newcommand{\code}[1]{\texttt{\small #1}}

% ---- Hilfsbefehl: zentrierte Folien-Abbildung ---------------
% #1 = Breite (Faktor von \linewidth), #2 = Dateiname, #3 = Caption
\newcommand{\folie}[3]{%
  \begin{figure}[h!]\centering
    \includegraphics[width=#1\linewidth]{#2}%
    \caption{#3}%
  \end{figure}}
